\section{班级情感建设}

此处开篇概述班级情感建设的理念，介绍本学年开展的主要情感建设活动。

\subsection{男生节 \& 女生节}

此处描述男生节当天的活动内容，如同学们准备的惊喜和庆祝横幅等。

% \begin{figure}[H]
%     \centering
%     \includegraphics[width=0.3\textwidth]{assets/example.jpg}
%     \includegraphics[width=0.3\textwidth]{assets/example.jpg}
%     \caption{男生节活动。左：横幅；右：班主任打卡祝福}
% \end{figure}

% \begin{figure}[H]
%     \centering
%     \includegraphics[width=0.3\textwidth]{assets/example.jpg}
%     \caption{男生节礼物}
% \end{figure}

此处描述女生节当天的活动内容和庆祝聚餐。

% \begin{figure}[H]
%     \centering
%     \includegraphics[width=0.45\textwidth]{assets/example.jpg}
%     \includegraphics[width=0.45\textwidth]{assets/example.jpg}
%     \caption{女生节活动。左：聚餐；右：横幅}
% \end{figure}

\subsection{日常交流}

此处介绍班级日常情感交流的举措，如在班主任支持下建立的交流机制（如约饭基金）。

此处描述活动的参与人次和对班级凝聚力的影响。

% \begin{figure}[H]
%     \centering
%     \includegraphics[width=0.3\textwidth]{assets/example.jpg}
%     \includegraphics[width=0.3\textwidth]{assets/example.jpg}
%     \includegraphics[width=0.3\textwidth]{assets/example.jpg}
%     \caption{班级活动}
% \end{figure}

\subsection{总结}

此处总结班级情感建设的意义和成效。
